\documentclass[italian,11pt,a4paper]{article}
\usepackage[T1]{fontenc}
\usepackage[utf8]{inputenc}
\usepackage{babel}
\usepackage[margin=2.5cm]{geometry}
\usepackage{array,tabularx}
\usepackage{float}
\usepackage{amssymb}
\usepackage{graphicx}
\usepackage[colorlinks=true,linkcolor=blue]{hyperref}

% Scheda vuota per un requisito funzionale. Duplicare e rinumerare se necessario.
\newcommand{\schedaRequisito}[1]{
\subsubsection{RE#1 --}
\begin{table}[H]\centering
\begin{tabularx}{\textwidth}{|>{\raggedright\arraybackslash}p{4cm}|X|}\hline
\multicolumn{2}{|c|}{\textbf{Utenti per funzionalità}}\\\hline
\textbf{Utente autenticato} & \textbf{Utente non autenticato}\\\hline
 & \\\hline
\end{tabularx}
\end{table}
\begin{table}[H]\centering
\begin{tabularx}{\textwidth}{|>{\raggedright\arraybackslash}p{4cm}|X|}\hline
Funzione & \\\hline
Descrizione & \\\hline
Input & \\\hline
Fonte & \\\hline
Output & \\\hline
Azione & \\\hline
Precondizione & \\\hline
Effetti collaterali & \\\hline
Use case & \\\hline
\end{tabularx}
\end{table}
}

\title{Specifica dei Requisiti}
\author{}
\date{}
\begin{document}
\maketitle
\tableofcontents
\newpage

\section{Introduzione}
\subsection{Scopo}

\subsection{Campo di applicazione}

\subsection{Definizioni (Glossario)}

\subsection{Riferimenti}

\newpage
\section{Descrizione generale}
\subsection{Prospettiva del prodotto}

\subsection{Funzionalità del prodotto}

\subsection{Attori del sistema}
\begin{table}[H]\centering
\begin{tabularx}{\textwidth}{|p{3.5cm}|X|X|}\hline
\multicolumn{3}{|c|}{\textbf{}}\\\hline
\textbf{Descrizione} & \textbf{Esperienza e competenze tecniche} & \textbf{Principali interazioni con il sistema}\\\hline
 & & \\\hline
\end{tabularx}
\end{table}

\begin{table}[H]\centering
\begin{tabularx}{\textwidth}{|p{3.5cm}|X|X|}\hline
\multicolumn{3}{|c|}{\textbf{}}\\\hline
\textbf{Descrizione} & \textbf{Esperienza e competenze tecniche} & \textbf{Principali interazioni con il sistema}\\\hline
 & & \\\hline
\end{tabularx}
\end{table}

\newpage
\section{Requisiti specifici}
\subsection{Requisiti funzionali}
\schedaRequisito{01}
\schedaRequisito{02}
\schedaRequisito{03}
\schedaRequisito{04}
\schedaRequisito{05}
\schedaRequisito{06}
\schedaRequisito{07}
\schedaRequisito{08}
\schedaRequisito{09}
\schedaRequisito{10}
\schedaRequisito{11}
\schedaRequisito{12}
\schedaRequisito{13}
\schedaRequisito{14}
\schedaRequisito{15}

\newpage
\subsection{Modellazione dei casi d'uso tramite Use Case Diagram}
\subsubsection{Use Case Diagram}

\subsubsection{Use Case Diagram parziale 1}
\begin{itemize}
\item
\end{itemize}
\begin{figure}[H]\centering
\fbox{\parbox[c][3cm][c]{0.75\textwidth}{\centering}}
\caption{}
\end{figure}

\subsubsection{Use Case Diagram parziale 2}
\begin{itemize}
\item
\end{itemize}
\begin{figure}[H]\centering
\fbox{\parbox[c][3cm][c]{0.75\textwidth}{\centering}}
\caption{}
\end{figure}

\subsubsection{Use Case Diagram parziale 3}
\begin{itemize}
\item
\end{itemize}
\begin{figure}[H]\centering
\fbox{\parbox[c][3cm][c]{0.75\textwidth}{\centering}}
\caption{}
\end{figure}

\subsubsection{Use Case Diagram parziale 4}
\begin{itemize}
\item
\end{itemize}
\begin{figure}[H]\centering
\fbox{\parbox[c][3cm][c]{0.75\textwidth}{\centering}}
\caption{}
\end{figure}

\subsubsection{Use Case Diagram parziale 5}
\begin{itemize}
\item
\end{itemize}
\begin{figure}[H]\centering
\fbox{\parbox[c][3cm][c]{0.75\textwidth}{\centering}}
\caption{}
\end{figure}

\subsubsection{Use Case Diagram parziale 6}
\begin{itemize}
\item
\end{itemize}
\begin{figure}[H]\centering
\fbox{\parbox[c][3cm][c]{0.75\textwidth}{\centering}}
\caption{}
\end{figure}

\newpage
\subsection{Requisiti non funzionali}
\subsubsection{RNF01 -- Sicurezza}

\subsubsection{RNF02 -- Affidabilità}

\subsubsection{RNF03 -- Prestazioni}

\subsubsection{RNF04 -- Usabilità}

\subsubsection{RNF05 -- Scalabilità}

\subsubsection{RNF06 -- Portabilità e indipendenza tra componenti}

\subsubsection{RNF07 -- Manutenibilità}

\newpage
\subsection{Formalizzazione di casi d'uso significativi}
\subsubsection{Tabella di Cockburn 1}
\begin{table}[H]\centering
\begin{tabularx}{\textwidth}{|p{4cm}|X|}\hline
\textbf{Use case} & \\\hline
Goal in context & \\\hline
Preconditions & \\\hline
Success end condition & \\\hline
Failed end condition & \\\hline
Trigger & \\\hline
Main scenario & \\\hline
Extensions & \\\hline
\end{tabularx}
\end{table}
\begin{table}[H]\centering
\begin{tabularx}{\textwidth}{|p{1.5cm}|X|X|}\hline
\textbf{Step} & \textbf{Attore} & \textbf{Risposta del sistema}\\\hline
 & & \\\hline
 & & \\\hline
 & & \\\hline
\end{tabularx}
\end{table}

\subsubsection{Tabella di Cockburn 2}
\begin{table}[H]\centering
\begin{tabularx}{\textwidth}{|p{4cm}|X|}\hline
\textbf{Use case} & \\\hline
Goal in context & \\\hline
Preconditions & \\\hline
Success end condition & \\\hline
Failed end condition & \\\hline
Trigger & \\\hline
Main scenario & \\\hline
Extensions & \\\hline
\end{tabularx}
\end{table}
\begin{table}[H]\centering
\begin{tabularx}{\textwidth}{|p{1.5cm}|X|X|}\hline
\textbf{Step} & \textbf{Attore} & \textbf{Risposta del sistema}\\\hline
 & & \\\hline
 & & \\\hline
 & & \\\hline
\end{tabularx}
\end{table}

\end{document}
